\section{Recommendations}
\label{sec:recs}

Ordered by effect on the legal position per unit of engineering effort. Items 1--5 are done and are
listed so that the record is complete; 6--12 are open.

\subsection*{Completed since the assessment}

\begin{enumerate}
  \item \textbf{Implement secure aggregation} and verify it end-to-end. Converts the answer to the
        central question from ``the server sees individual updates'' into a configuration choice.
        (\S\ref{sec:secagg})
  \item \textbf{Implement and measure local DP.} Establishes what noise-at-the-practice actually
        costs rather than asserting availability. The answer is uncomfortable and is reported as
        such. (\S\ref{sec:noise})
  \item \textbf{Replace the privacy accounting} with a standard R\'enyi-DP accountant. The previous
        figures were invalid at low noise, not merely loose. (\S\ref{sec:budget})
  \item \textbf{Build the membership-inference audit}, with a positive control. Converts ``DP is
        configured'' into a reported attack result. (\S\ref{sec:anontest})
  \item \textbf{Close the two side channels} --- the direct identifier in the export, and identified
        per-practice metric logging, now anonymised by default. (\S\ref{sec:sidechannels})
\end{enumerate}

\subsection*{Open, in priority order}

\begin{enumerate}[start=6]
  \item \textbf{Decide who operates the SuperLink, and decide secure aggregation with it.} These are
        the same decision viewed from two sides, they gate the Art.~26/28 analysis, and they have the
        longest lead time of anything on this list. (\S\ref{sec:outside})
  \item \textbf{Build the remaining two attack families} --- model inversion and reconstruction ---
        so that the audit covers what EDPB Opinion 28/2024 actually names. Consider also a
        likelihood-ratio attack as a stronger membership test. (\S\ref{sec:anontest})
  \item \textbf{Obtain the analytics work package's indicator specification} and apply the same
        accountant and a cell-suppression analysis to it. Until then, half of the assessment's
        framing is unanswered. (\S\ref{sec:aggdisclosure})
  \item \textbf{Confirm the round count on real data.} Convergence by round~10 is the cheapest
        privacy improvement available and it is currently a synthetic-data finding.
        (\S\ref{sec:reduce})
  \item \textbf{Evaluate subsampling amplification.} Participation below 1.0 earns a real budget
        reduction the accountant can now credit; the cost is convergence speed. (\S\ref{sec:reduce})
  \item \textbf{Re-run the whole evidence base on the pilot cohort} once real data exists ---
        accuracy, budget and audit alike. The right-hand column of \S\ref{sec:realdata} is the
        checklist.
  \item \textbf{Implement FedMosaic's own DP construction} if that protocol is to be considered on
        privacy grounds, since its current case rests on payload shape alone. (\S\ref{sec:payload})
\end{enumerate}

\subsection*{Mapping to the project plan}

\begin{center}
\small
\begin{tabular}{lll}
  \toprule
  Item & Work package & Milestone \\
  \midrule
  Secure aggregation (1), operator decision (6) & T2.3 & MS4 (month 18) \\
  DP measurement (2), accounting (3), amplification (11) & T2.3 & MS4 \\
  Leakage audit (4), remaining attacks (7) & T2.5 & MS4 onward \\
  Side channels (5) & T2.3 / data collection & before MS3 \\
  Analytics indicator spec (8) & analytics work package & --- \\
  Re-run on pilot cohort (12) & T2.4 & MS5 (month 24) \\
  \bottomrule
\end{tabular}
\end{center}
